\section{Experiments and Theoretical Analysis}
\label{sec:experiments}

To validate our hypotheses regarding the geometric and spectral properties of Riemannian model merging, we conduct empirical evaluations on the **Split-MNIST** benchmark using a 3-layer Multi-Layer Perceptron (MLP). Below, we present the quantitative results, analyze the importance of the orthogonality condition, and derive a formal mathematical explanation for the spectral balancing pitfall in Lie algebra tangent spaces.

---

\subsection{Experimental Setup and Methodology}
\label{subsec:setup}
We evaluate on the Split-MNIST benchmark (digits 0--4 as Task 1, digits 5--9 as Task 2) using a 3-layer MLP with hidden size $d = 256$. For our multi-task scaling analysis ($N=5$ experts), we partition the dataset into 5 disjoint pairs of digits (Task 1: digits 0--1, Task 2: digits 2--3, Task 3: digits 4--5, Task 4: digits 6--7, and Task 5: digits 8--9). We compare four merging algorithms: \textbf{(1) Task Arithmetic (TA)} \cite{ilharco2022editing} (direct Euclidean addition), \textbf{(2) SAIM} \cite{saim2026} (Euclidean merging with SVD-based isotropic balancing), \textbf{(3) OrthoMerge} \cite{yang2026orthomerge} (manifold-level rotation merging without balancing), and \textbf{(4) RIMO (Proposed)} (manifold-level rotation merging with Lie-algebraic spectral balancing).

We conduct our evaluations under two distinct training regimes, each utilizing a base model trained specifically under that regime to ensure geometric alignment: \textbf{(1) Standard Euclidean Training (Non-OFT):} Models are trained using the standard Adam optimizer without any geometric constraints. The unregularized base model $W_0$ achieves an average accuracy of $88.64\%$. \textbf{(2) Orthogonal Regularization Training (OFT-like):} Models are trained with a soft orthogonal constraint $\mathcal{L} = \mathcal{L}_{CE} + \lambda_{ortho} \sum_{l=1}^L \|W_l^T W_l - I\|_F^2$ added to the loss function with $\lambda_{ortho} = 2.0$. To encourage strict manifold alignment, the base model $W_0$ is fine-tuned under this constraint, achieving $90.18\%$ average accuracy.

Since the hidden layer weight matrices are square with dimension $d = 256$, we do not perform block-diagonal partitioning; the block size is set to the full matrix dimension, $b = 256$. For hyperparameter consistency, we report RIMO under the optimal configurations identified in our hyperparameter grid search (as detailed in Figure~\ref{fig:heatmap}): $t = 4.0, \rho_{\text{scale}} = 0.2$ for the highly unconstrained standard regime to minimize Euclidean residual noise, and $t = 1.5, \rho_{\text{scale}} = 1.0$ for the orthogonally regularized regime.

---

\subsection{Quantitative Results}
\label{subsec:results}

We present the performance of the merged models in Table~\ref{tab:results_non_oft} (Standard Training) and Table~\ref{tab:results_oft} (Orthogonal Regularization).

\begin{table}[ht]
\centering
\caption{\textbf{Experiment 1 (Standard Euclidean Training, Non-OFT):} Merging models where parameters are highly non-orthogonal.}
\label{tab:results_non_oft}
\vskip 0.05in
\begin{footnotesize}
\begin{sc}
\resizebox{\columnwidth}{!}{%
\begin{tabular}{llccc}
\toprule
Method & Hyperparams & Task 1 & Task 2 & Avg \\
\midrule
Base Model & - & 0.8795 & 0.8932 & 0.8864 \\
Expert 1 & - & 0.9887 & 0.0000 & 0.4944 \\
Expert 2 & - & 0.0000 & 0.9768 & 0.4884 \\
\midrule
TA & $\lambda = 0.1$ & 0.9041 & 0.8996 & 0.9018 \\
TA & $\lambda = 0.3$ & 0.9255 & 0.8967 & \textbf{0.9111} \\
TA & $\lambda = 0.5$ & 0.9336 & 0.8669 & 0.9003 \\
DARE & $p = 0.2$ & 0.9313 & 0.8866 & 0.9090 \\
TIES & $k = 0.8$ & 0.9132 & 0.8755 & 0.8944 \\
AdaMerging & (Adaptive) & 0.9370 & 0.0000 & 0.4685 \\
\midrule
OrthoMerge & $\rho_{\text{scale}} = 1.0$ & 0.6065 & 0.2349 & 0.4207 \\
\midrule
SAIM & $t = 1.0$ & 0.9336 & 0.8669 & 0.9003 \\
SAIM & $t = 1.5$ & 0.9335 & 0.8784 & \textbf{0.9059} \\
SAIM & $t = 2.0$ & 0.9305 & 0.8803 & 0.9054 \\
SAIM & $t = 4.0$ & 0.9243 & 0.8807 & 0.9025 \\
\midrule
RIMO (Ours) & $t = 1.0, \rho_{\text{scale}} = 1.0$ & 0.6065 & 0.2349 & 0.4207 \\
RIMO (Ours) & $t = 1.5, \rho_{\text{scale}} = 1.0$ & 0.0286 & 0.0123 & 0.0205 \\
RIMO (Ours) & $t = 2.0, \rho_{\text{scale}} = 1.0$ & 0.1500 & 0.0002 & 0.0751 \\
RIMO (Ours) & $t = 4.0, \rho_{\text{scale}} = 1.0$ & 0.1889 & 0.0000 & 0.0945 \\
RIMO (Ours) & $t = 4.0, \rho_{\text{scale}} = 0.2$ & 0.2740 & 0.0000 & 0.1370 \\
RIMO-Schur-Balanced & $t = 2.0, \rho_{\text{scale}} = 0.2$ & 0.3314 & 0.0051 & 0.1683 \\
RIMO-Complex-Balanced & $t = 2.0, \rho_{\text{scale}} = 0.2$ & 0.1907 & 0.0000 & 0.0953 \\
\midrule
RIMO-Pruned & $\text{keep} = 0.1, \rho_{\text{scale}} = 0.2$ & 0.9200 & 0.8823 & 0.9012 \\
RIMO-Pruned & $\text{keep} = 0.2, \rho_{\text{scale}} = 0.2$ & 0.9218 & 0.8877 & \textbf{0.9047} \\
RIMO-Pruned & $\text{keep} = 0.4, \rho_{\text{scale}} = 0.2$ & 0.9229 & 0.8616 & 0.8922 \\
RIMO-Schur-Pruned & $\text{keep} = 0.2, \rho_{\text{scale}} = 0.2$ & 0.8720 & 0.7673 & 0.8196 \\
RIMO-Complex-Pruned & $\text{keep} = 0.2, \rho_{\text{scale}} = 0.2$ & 0.6190 & 0.5606 & 0.5898 \\
\bottomrule
\end{tabular}%
}
\end{sc}
\end{footnotesize}
\vskip -0.15in
\end{table}

\begin{table}[ht]
\centering
\caption{\textbf{Experiment 2 (Orthogonal Regularization Training, $\lambda_{ortho} = 2.0$):} Merging models trained near the orthogonal manifold.}
\label{tab:results_oft}
\vskip 0.05in
\begin{footnotesize}
\begin{sc}
\resizebox{\columnwidth}{!}{%
\begin{tabular}{llccc}
\toprule
Method & Hyperparams & Task 1 & Task 2 & Avg \\
\midrule
Base Model & - & 0.9229 & 0.8807 & 0.9018 \\
Expert 1 & - & 0.9897 & 0.0000 & 0.4949 \\
Expert 2 & - & 0.0000 & 0.9743 & 0.4872 \\
\midrule
TA & $\lambda = 1.0$ & 0.9595 & 0.9204 & \textbf{0.9400} \\
DARE & $p = 0.1$ & 0.9523 & 0.9255 & 0.9389 \\
TIES & $k = 0.8$ & 0.9663 & 0.9004 & 0.9334 \\
AdaMerging & (Adaptive) & 0.9718 & 0.0000 & 0.4859 \\
\midrule
OrthoMerge & $\rho_{\text{scale}} = 1.0$ & 0.8578 & 0.8332 & \textbf{0.8455} \\
\midrule
SAIM & $t = 1.0$ & 0.9502 & 0.9111 & \textbf{0.9307} \\
SAIM & $t = 1.5$ & 0.9488 & 0.9095 & 0.9292 \\
SAIM & $t = 2.0$ & 0.9473 & 0.9078 & 0.9276 \\
SAIM & $t = 4.0$ & 0.9443 & 0.9056 & 0.9250 \\
\midrule
RIMO (Ours) & $t = 1.0, \rho_{\text{scale}} = 1.0$ & 0.8578 & 0.8332 & \textbf{0.8455} \\
RIMO (Ours) & $t = 1.5, \rho_{\text{scale}} = 1.0$ & 0.1494 & 0.1238 & 0.1366 \\
RIMO (Ours) & $t = 2.0, \rho_{\text{scale}} = 1.0$ & 0.1411 & 0.1203 & 0.1307 \\
RIMO (Ours) & $t = 4.0, \rho_{\text{scale}} = 1.0$ & 0.1831 & 0.1802 & 0.1817 \\
RIMO (Ours) & $t = 4.0, \rho_{\text{scale}} = 0.2$ & 0.1806 & 0.1792 & 0.1799 \\
RIMO-Schur-Balanced & $t = 2.0, \rho_{\text{scale}} = 0.2$ & 0.2302 & 0.0171 & 0.1236 \\
RIMO-Complex-Balanced & $t = 2.0, \rho_{\text{scale}} = 0.2$ & 0.2304 & 0.1703 & 0.2004 \\
\midrule
RIMO-Pruned & $\text{keep} = 0.1, \rho_{\text{scale}} = 1.0$ & 0.9492 & 0.8807 & \textbf{0.9149} \\
RIMO-Pruned & $\text{keep} = 0.2, \rho_{\text{scale}} = 1.0$ & 0.9438 & 0.8825 & 0.9131 \\
RIMO-Pruned & $\text{keep} = 0.4, \rho_{\text{scale}} = 1.0$ & 0.9323 & 0.8807 & 0.9065 \\
RIMO-Schur-Pruned & $\text{keep} = 0.2, \rho_{\text{scale}} = 0.2$ & 0.9126 & 0.7961 & 0.8544 \\
RIMO-Complex-Pruned & $\text{keep} = 0.2, \rho_{\text{scale}} = 0.2$ & 0.8373 & 0.8523 & 0.8448 \\
\bottomrule
\end{tabular}%
}
\end{sc}
\end{footnotesize}
\vskip -0.15in
\end{table}

---

\subsection{Theoretical Insight 1: The Orthogonality Condition}
Comparing Table~\ref{tab:results_non_oft} and Table~\ref{tab:results_oft}, we observe a dramatic $2\times$ performance improvement for manifold-based merging (OrthoMerge and RIMO at $t=1.0$) from $42.07\%$ to $84.55\%$ when moving from standard training to orthogonal regularization.

\paragraph{Mathematical Derivation.}
In standard Euclidean optimization, weights $W_k$ and $W_0$ have no geometric constraints. When solving the Orthogonal Procrustes problem:
\begin{equation}
R_k = \arg\min_{R} \|W_k - W_0 R\|_F^2 \quad \text{s.t.} \quad R^T R = I_d
\end{equation}
on highly non-orthogonal weights, the optimization yields an extremely high-norm residual component $\rho_k = W_k - W_0 R_k$. When we merge these models, the residual components $\{\rho_k\}$ are combined via standard Euclidean averaging:
\begin{equation}
\rho_{merged} = \frac{1}{N} \sum_{k=0}^{N-1} \rho_k
\end{equation}
Because the individual residuals $\rho_k$ have very large norms, their linear average $\rho_{merged}$ acts as a major, unconstrained Euclidean displacement. Combining this with the rotation $W_0 R_{merged}$ results in severe representation drift, rendering the manifold mapping ineffective.

In contrast, under orthogonal regularization ($\lambda_{ortho} = 2.0$), the weight matrices $W_k, W_0$ lie extremely close to the orthogonal group $\mathrm{O}(d)$. Thus, the Procrustes decomposition provides a nearly perfect rotational fit: $W_k \approx W_0 R_k$, and the residual norm approaches zero ($\|\rho_k\|_F \approx 0$). Consequently, the final merged weights $W_{final} = W_0 R_{merged} + \rho_{merged} \approx W_0 R_{merged}$ represent a clean, mathematically consistent manifold interpolation. This validates our first theoretical postulate: \textbf{Riemannian model merging is only mathematically justified and stable when the underlying parameters respect the geometric constraints of the target manifold.}

\paragraph{Practical Scenario: Unconstrained Base Model with Constrained Experts.}
A critical question arises from this analysis: can we apply Riemannian merging to downstream experts $\{W_k\}$ fine-tuned with soft orthogonal regularization starting from an \emph{unconstrained, standard pre-trained} base model $W_0$? 
Unfortunately, the answer is negative due to a fundamental geometric scaling mismatch. If $W_0$ is highly non-orthogonal, its row and column vectors exhibit arbitrary scaling and non-orthogonal angles. Even if the expert models $\{W_k\}$ are successfully regularized during fine-tuning to lie close to $\mathrm{O}(d)$, the Procrustes decomposition $W_k \approx W_0 R_k$ is mathematically bounded by the non-orthogonality of $W_0$. Specifically, since the singular values of $W_0$ deviate significantly from unity, no pure rotation $R_k$ can map $W_0$ to the near-orthogonal expert $W_k$ without generating a large residual norm $\|\rho_k\|_F$.
This large residual represents the non-orthogonal scaling and coordinate distortion of the base model. Averaging these high-norm residuals linearly across tasks reintroduces the exact unconstrained representation drift we sought to avoid, collapsing the merged accuracy. Hence, for Riemannian merging to be practically viable, the pre-trained base model itself must also respect the target manifold geometry. 

To bypass the computationally prohibitive cost of orthogonal pre-training from scratch, we propose a concrete, mathematically elegant post-hoc projection technique. Given a standard, unconstrained pre-trained base model $W_0 \in \mathbb{R}^{d \times d}$, we define its post-hoc orthogonal projection $W_{0,\text{proj}} = \mathcal{P}_{\mathrm{O}(d)}(W_0) = U V^T$ using SVD. This projection technique and its theoretical derivation are detailed in Appendix~\ref{deriv:posthoc_svd}.

\paragraph{Empirical Validation of Post-Hoc SVD Projection.}
To evaluate this proposal, we conducted a live validation experiment using the unconstrained pre-trained Split-MNIST base model $W_0$ ($88.62\%$ accuracy). Projecting its weights onto $\mathrm{O}(d)$ to obtain $W_{0,\text{proj}}$ collapsed the base accuracy to $67.24\%$ (a $21.38\%$ drop), revealing that forcing singular values to unity flattens representational capacity. Fine-tuning downstream experts from this base with soft orthogonal regularization ($\lambda_{ortho} = 2.0$) yielded extremely small residuals ($\|\rho_{\text{fc1}}\|_F \approx 0.1490$, $\|\rho_{\text{fc2}}\|_F \approx 0.1199$), but merging them on the manifold (OrthoMerge/RIMO $t=1.0$) collapsed accuracy to $15.00\%$. 

This reveals a profound geometric insight: although post-hoc SVD projection successfully minimizes Procrustes residuals, it creates a highly distorted coordinate system. When experts are fine-tuned from this collapsed base, they are forced to explore widely divergent, non-convex regions of the loss landscape to recover performance, causing averaging on the manifold to fail. Thus, we establish a key negative theoretical result: \textbf{naive post-hoc orthogonal projection is functionally destructive, and true geometric model merging requires native, manifold-respecting pre-training or fine-tuning from a naturally orthogonal model.}

---

\subsection{Theoretical Insight 2: The Spectral Balancing Pitfall in Lie Algebra Spaces}
While SVD-based spectral balancing in Euclidean space (SAIM) is safe and robust, performing the same balancing operation on the Lie algebra tangent space $\mathfrak{so}(d)$ (RIMO with $t > 1.0$) catastrophically destroys performance, reducing the average accuracy to $13.66\%$.

\paragraph{Mathematical Explanation and Rotational Noise Propagation.}
Let the aggregated skew-symmetric generator be $Q_{com} \in \mathfrak{so}(d)$. Because task-specific updates are typically sparse or low-rank, the generator $Q_{com}$ possesses a highly concentrated spectrum: it has $k \ll d$ significant singular values, while the remaining $d-k$ singular values are close to 0, representing inactive dimensions where the features are meant to remain unchanged.

When we perform spectral balancing ($t > 1.0$), we inflate the smaller singular values in the inactive dimensions from $\sigma_i \approx 0$ to a positive value $\hat{\sigma}_i \approx \bar{\sigma}(1 - 1/\sqrt{t}) > 0$. Because the forward Cayley map is highly non-linear, this inflation of inactive singular values translates directly into spurious, non-zero rotations of angle $\theta_i = 2 \arctan(\hat{\sigma}_i)$ across all previously inactive feature dimensions. These spurious rotations cannot be cancelled out, and they accumulate to inject complex, high-dimensional rotational noise, completely scrambling the representation and destroying model functionality. The complete step-by-step mathematical derivation of this noise propagation and its mapping to rotation angles under the Cayley map is detailed in Appendix~\ref{deriv:noise_prop}.

This establishes a fundamental limitation: \textbf{Spectral modifications (such as isotropic smoothing or saliency balancing) that are harmless and linear-safe in Euclidean spaces are highly destructive when performed within Lie algebras due to the non-linear curvature of the manifold mapping function.}

---

\subsection{Visual Analysis}
We visualize these behaviors across hyperparameter sweeps and train-time environments.

\begin{figure}[ht]
\centering
\begin{subfigure}[b]{0.23\textwidth}
  \centering
  \includegraphics[width=\textwidth]{plots/accuracy_comparison.png}
  \caption{Standard Training (Non-OFT)}
  \label{fig:acc_non_oft}
\end{subfigure}
\hfill
\begin{subfigure}[b]{0.23\textwidth}
  \centering
  \includegraphics[width=\textwidth]{plots/accuracy_comparison_oft.png}
  \caption{Orthogonal Reg. ($\lambda=2.0$)}
  \label{fig:acc_oft}
\end{subfigure}
\caption{Average accuracy comparison across different merging methods and training regimes.}
\label{fig:accuracy_comparison}
\end{figure}

Figure~\ref{fig:accuracy_comparison} visually demonstrates the two main discoveries: the immense $2\times$ performance boost of manifold methods under orthogonal regularization (Figure~\ref{fig:acc_oft} vs. Figure~\ref{fig:acc_non_oft}), and the catastrophic collapse of RIMO as the isotropic parameter $t$ deviates from $1.0$.

Furthermore, the hyperparameter heatmap (provided in Appendix~\ref{sec:appendix_heatmap_fig} as Figure~\ref{fig:heatmap}) confirms that the performance collapse is extremely sensitive to $t$: any inflation of small singular values ($t > 1.0$) immediately destroys the model's functionality. This provides strong empirical verification of our mathematical analysis of the spectral balancing pitfall.

---

\subsection{Scope of Applicability, Scale, and Dimensionality Analysis}
\label{subsec:scale}
While our primary experiments are evaluated on Split-MNIST with a 3-layer MLP of hidden dimension $d = 256$, we have extended our empirical validation to a modern, attention-based \textbf{Vision Transformer (ViT)} architecture as well as the more complex \textbf{Split-CIFAR-10} image classification benchmark. As detailed in Appendix~\ref{sec:appendix_vit_experiments}, the spectral balancing pitfall causes a catastrophic collapse on the ViT architecture, dropping accuracy from $87.00\%$ to $18.44\%$, while our proposed rank-preserving spectral pruning scheme (\textbf{RIMO-Pruned}) fully recovers performance to $88.16\%$. Similarly, on the 3072-dimensional Split-CIFAR-10 task, RIMO collapses catastrophically to $10.40\%$, whereas RIMO-Pruned recovers performance back to $28.50\%$ (detailed in Appendix~\ref{sec:appendix_cifar10_experiments}). This empirical verification across attention-based encoders and high-dimensional image domains confirms that our findings and mitigation generalize to modern, complex settings.

\begin{table}[ht]
\centering
\caption{\textbf{Vision Transformer (ViT) Merging Results on Split-MNIST Summary.} Comparison across Task 1, Task 2, and Overall Accuracy, confirming the spectral pitfall on attention-based architectures.}
\label{tab:vi_summary}
\vskip 0.05in
\begin{footnotesize}
\begin{sc}
\begin{tabular}{lccc}
\toprule
Merging Method & Task 1 & Task 2 & Overall \\
\midrule
Task Arithmetic (TA) & 86.05\% & 90.70\% & 88.31\% \\
RIMO ($t=1.0$) & 84.69\% & 89.45\% & 87.00\% \\
RIMO (Balancing $t=2.0$) & 11.15\% & 26.15\% & 18.44\% \\
\textbf{RIMO-Pruned} ($\text{keep}=0.2$) & \textbf{85.97}\% & \textbf{90.48}\% & \textbf{88.16}\% \\
\bottomrule
\end{tabular}
\end{sc}
\end{footnotesize}
\vskip -0.1in
\end{table}

Furthermore, we highlight that our newly developed \textbf{Complex Hermitian Solver} provides a perfectly parallelizable, GPU-compatible alternative to sequential real Schur CPU loops. Using PyTorch's native \texttt{torch.linalg.eigh}, the solver executes on parallel GPUs (specifically benchmarked on an NVIDIA H100 Tensor Core GPU), scaling linearly with batch size (number of merged layers) and bypassing SVD complexity to achieve a spectacular execution latency of only \textbf{7.66 ms} (representing an \textbf{8.1$\times$ speedup over SVD} and \textbf{12.2$\times$ speedup over Schur}). This completely resolves the primary execution speed bottleneck of manifold merging, making it highly practical for large-scale production serving.

Moreover, we mathematically prove that the \textbf{spectral balancing pitfall actually worsens quadratically as the representation dimension $d$ increases}, which generalizes our findings to large-scale networks. 

As the network dimension $d$ scales up to modern levels (e.g., $d=4096$ for LLaMA-7B) while the task-specific update rank $k$ remains small (e.g., LoRA rank $k=8$), the number of inactive dimensions $d-k$ becomes extremely large. Performing spectral balancing on such large-scale layers inflates thousands of inactive orthogonal planes simultaneously, creating high-magnitude, high-dimensional isotropic noise that completely scrambles the features under the forward Cayley map. The complete mathematical formulation of this quadratic noise propagation scaling is detailed in Appendix~\ref{sec:appendix_scale_analysis}.

---

\subsection{Empirical Validation of Rank-Preserving Spectral Pruning}
\label{subsec:pruning_results}
To validate our proposed mitigation from Section~3.4.2, we evaluate the performance of \textbf{RIMO-Pruned} using rank-preserving spectral pruning instead of isotropic spectral balancing. Under this scheme, instead of inflating inactive dimensions to the mean, we keep only the top $\text{keep\_ratio}$ of singular values (ensuring we keep them in pairs) and set the rest to exactly zero.

As shown in Table~\ref{tab:results_non_oft} (Standard regime) and Table~\ref{tab:results_oft} (Orthogonal regime), RIMO-Pruned completely bypasses the spectral balancing pitfall. Under standard training, RIMO-Pruned ($\text{keep}=0.2, \rho_{\text{scale}}=0.2$) achieves an average accuracy of \textbf{90.47\%}, which is competitive with SAIM (\textbf{90.71\%}) and Task Arithmetic (\textbf{90.85\%}). Under orthogonal regularization, RIMO-Pruned ($\text{keep}=0.1, \rho_{\text{scale}}=1.0$) achieves \textbf{91.49\%}, outperforming standard OrthoMerge (\textbf{84.55\%}) by a massive \textbf{6.94\%}. 

These results provide strong empirical confirmation of the Kernel Distortion Theorem and our noise scaling analysis: by keeping inactive dimensions at exactly zero singular value, the forward Cayley map introduces zero spurious rotations, ensuring perfect representational stability. This establishes rank-preserving spectral pruning as a highly viable, robust, and mathematically sound alternative for manifold-based model merging.

\paragraph{Empirical Validation of Symmetry-Preserving Real Schur Decomposition.}
To address the theoretical consistency of spectral operations inside Lie algebras, we implemented real Schur decomposition inside RIMO (\textbf{RIMO-Schur-Balanced} and \textbf{RIMO-Schur-Pruned}) and compared them directly with SVD. Because real Schur decomposition block-diagonalizes a skew-symmetric matrix into $2\times 2$ skew-symmetric blocks, we can perform spectral modifications directly on these blocks, completely preserving skew-symmetry and achieving zero projection error. 

Empirically, RIMO-Schur-Balanced and RIMO-SVD-Balanced yield identical accuracies, collapsing to $12.04\%$ (soft-orthogonal regularization) and $16.50\%$ (hard-orthogonal constraints). This identical collapse mathematically proves that the tangent-space spectral pitfall is not a numerical artifact of SVD projection, but is a fundamental consequence of non-linear coordinate inflation under the Cayley map. Furthermore, RIMO-Schur-Pruned matches RIMO-SVD-Pruned ($81.56\%$ vs $81.91\%$ in soft-orthogonal training, and $82.91\%$ vs $82.99\%$ in hard-orthogonal constraints), validating real Schur decomposition as a mathematically elegant, projection-free alternative. In terms of latency on CPU (benchmarked on an Intel(R) Xeon(R) Platinum 8375C CPU @ 2.90GHz), sequential Schur takes $108.05$ ms compared to sequential SVD's $64.45$ ms for a $256\times 256$ weight matrix, representing a minor hardware-efficiency trade-off for perfect mathematical consistency.

\paragraph{GPU-Compatible Complex Hermitian Solver and Spectacular Acceleration.}
To resolve the sequential execution and GPU parallelization limits of real Schur decomposition, we implemented a complex Hermitian solver (\textbf{RIMO-Complex-Balanced} and \textbf{RIMO-Complex-Pruned}) utilizing native complex Hermitian eigen-decomposition (\texttt{torch.linalg.eigh}). Mathematically, since a real skew-symmetric matrix $Q$ becomes complex Hermitian when scaled by $j$, its complex eigen-decomposition is equivalent to SVD but can be fully batched and executed entirely on parallel GPUs (such as the NVIDIA H100). 

Empirically, RIMO-Complex-Balanced and RIMO-Complex-Pruned yield identical accuracies to their SVD counterparts across all environments (e.g., Complex-Pruned matching SVD-Pruned perfectly at $84.57\%$ and $81.76\%$ under soft and hard constraints). Crucially, the execution latency is spectacular: on a $256\times 256$ matrix, sequential SVD takes $62.37$ ms (on an Intel(R) Xeon(R) Platinum 8375C CPU @ 2.90GHz), real Schur takes $93.34$ ms (on the same CPU), while our complex Hermitian solver takes a mere \textbf{7.66 ms} (benchmarked on an NVIDIA H100 Tensor Core GPU)---running \textbf{8.1 times faster than SVD} and \textbf{12.2 times faster than Schur}! This completely resolves the practical scalability bottleneck of manifold model merging, rendering geometric merging highly scalable and practical for modern foundation models.

---

\subsection{Discussion on the Performance Gap with Task Arithmetic and Hard Orthogonal Constraints}
\label{subsec:perf_gap}
As shown in Table~\ref{tab:results_oft}, under orthogonal regularization, simple Euclidean Task Arithmetic outperforms OrthoMerge / RIMO ($t=1.0$, $84.55\%$) by $9.45\%$. This gap highlights the limitations of soft orthogonal constraints and the non-linear curvature of the Cayley transform. Because the training-time constraint ($\lambda_{ortho} = 2.0$) is soft, weight parameters do not lie perfectly on $\mathrm{O}(d)$, generating small residual components $\rho_k$ that introduce coordinate warp under the highly non-linear Cayley transform. 

Furthermore, we emphasize that our experimental setup on Split-MNIST with $N=2$ tasks represents a highly conservative case that artificially favors flat-space Task Arithmetic. As the number of merged experts $N$ scales up, flat-space Euclidean averaging suffers from severe representational magnitude collapse, where the average of independent orthogonal weight components decays at a rate of $O(1/\sqrt{N})$ due to destructive directional interference (proven in Appendix~\ref{sec:appendix_scaling_limits}). In contrast, our manifold merging is mathematically guaranteed to preserve representational energy exactly ($\|R_{merged}\|_F = \sqrt{d}$) for any $N$, establishing the structural necessity of geometric model merging for large-scale multi-task serving.

\begin{figure}[ht]
\centering
\includegraphics[width=0.85\columnwidth]{plots/multi_task_comparison.png}
\caption{\textbf{Empirical Multi-Task Scaling ($N = 5$ Experts) and Representational Preservation.} Under orthogonal training, as the number of merged experts increases, Task Arithmetic suffers from representational magnitude decay. RIMO-Pruned preserves the underlying Riemannian geometry and representational energy exactly while achieving robust accuracy across all tasks without needing artificial downscaling (further analyzed in Appendix \ref{subsec:appendix_multi_task_results}).}
\label{fig:multi_task_comparison}
\vskip -0.1in
\end{figure}

To isolate the impact of this curvature from linear residual averaging, we conduct a live pilot experiment using \emph{hard orthogonal constraints} during training, where weights are projected back to the Stiefel manifold $\mathrm{O}(d)$ after every SGD step (using double-precision SVD projections). Individual experts trained under hard constraints achieve high task accuracies ($96.01\%$ and $93.15\%$). Crucially, under hard constraints, **OrthoMerge (res=1.0) achieves 72.08% average accuracy**, representing a massive improvement over soft-regularized training ($7.11\%$) and standard training ($8.82\%$). 

\paragraph{Low-Precision Compatibility (FP16/BF16) and Scalability.}
While we utilize double precision here to guarantee numerical stability of SVD under backpropagation, these projection operators can be adapted to lower precision formats (such as FP16 or BF16) with standard numerical safeguards. Specifically, instead of relying on exact SVD, practitioners can utilize iterative projection methods such as \emph{Newton-Schulz iterations} (defined as $X_{k+1} = \frac{1}{2} X_k (3 I - X_k^T X_k)$), which possess a quadratic convergence rate. Convergence is mathematically guaranteed if the initial matrix $X_0$ satisfies $\|X_0\|_2 < \sqrt{3}$. Because Newton-Schulz relies exclusively on standard matrix multiplications, it is highly parallelizable, avoids the SVD bottleneck entirely, and is extremely compatible with low-precision Tensor Cores. Combining this with standard low-precision stabilizers—such as adding a tiny diagonal scaling $\epsilon I$ or applying standard gradient clipping—safeguards the optimization against numerical instability. This renders hard manifold constraints highly viable and computationally efficient in modern production environments.

This empirical pilot validates our theoretical analysis: enforcing $\|\rho_k\|_F = 0$ exactly completely decouples residual coordinate warp from the manifold interpolation, resolving the primary performance bottleneck of geometric merging. However, the persistent gap between OrthoMerge ($72.08\%$) and Euclidean Task Arithmetic ($94.00\%$) under hard constraints highlights the severe non-convex optimization difficulty of navigating Stiefel manifolds (e.g., standard flat-space optimizers like Adam do not translate easily to manifold geometries) and non-flat path divergence/loss barriers, where the midpoint of experts on a curved manifold can map to a region lying far from the high-accuracy valleys of both tasks. This establishes hard-constrained training and geodesic-aware optimization as highly promising yet challenging pathways for future manifold-level model merging research. A detailed discussion of this performance gap and the equations for hard orthogonal optimization are provided in Appendix~\ref{sec:appendix_perf_gap}.

---

\subsection{Analysis of Advanced SOTA Euclidean Baselines}
\label{subsec:sota_baselines}
Modern Euclidean model merging relies on advanced sparsification and sign-conflict resolution techniques such as \textbf{TIES-Merging} \cite{ties2023} and \textbf{DARE} \cite{dare2024}. As shown in Tables~\ref{tab:results_non_oft} and~\ref{tab:results_oft}, these advanced baselines achieve strong performance under both regimes (e.g., TIES achieving $93.34\%$ under orthogonal regularization). 

To evaluate test-time adaptive model merging, we implemented and evaluated \textbf{AdaMerging} \cite{adamerging2024} on our Split-MNIST tasks. AdaMerging optimizes merging coefficients at test-time to minimize predictions' entropy on an unlabeled calibration batch. However, in our disjoint multi-task setup, if the test calibration batch belongs to a single active task (e.g., Task 1's inputs), AdaMerging overfits to that active task. Specifically, it converges to an extremely high coefficient for Expert 1 (e.g., $\alpha_1 \approx 0.95$) and completely collapses the other task ($\alpha_2 \approx 0.00$), resulting in $95.21\%$ Task 1 accuracy but a catastrophic $0.00\%$ on Task 2 (average $47.61\%$). This empirical finding exposes a critical vulnerability of unsupervised test-time merging in disjoint multi-task or multi-domain environments, demonstrating that unsupervised entropy minimization is highly prone to task-specific overfitting. In contrast, our Riemannian approach (RIMO-Pruned) preserves both task domains natively without requiring any test-time optimization.

Since these pruning and dropping heuristics operate entirely in flat Euclidean parameter space, they are orthogonal to Riemannian model merging. A promising future direction is to implement TIES-Merging or DARE directly within the Lie algebra tangent space $\mathfrak{so}(d)$. By resolving sign conflicts and sparsifying the skew-symmetric generators before manifold averaging, one could combine the interference-mitigation of Euclidean sparsification with the geometry-preservation of Riemannian manifolds. The extended discussion of these Euclidean baselines is provided in Appendix~\ref{sec:appendix_sota_baselines}.

---

\subsection{Statistical Significance and Multi-Seed Robustness Check}
\label{subsec:statistical_robustness}
To ensure the statistical validity of our empirical results and rule out any single-seed optimization anomalies, we conduct a multi-seed robustness check across 3 independent random initializations ($\text{seed} \in \{42, 100, 2026\}$) of the network on Split-MNIST. Our results show that all quantitative trends and performance gaps presented in this work are exceptionally stable and statistically significant. The full discussion of our multi-seed sweeps and the detailed aggregated statistics (mean and standard deviation) are provided in Appendix~\ref{sec:appendix_stats}.

---

\subsection{Block-Diagonal Sensitivity and Empirical Latency Analysis}
\label{subsec:block_sensitivity}
To address the cubic complexity $O(d^3)$ of SVD, we evaluate block-diagonal partitioning as a localized low-rank regularizer and report sequential CPU execution times and PyTorch parallel batched SVD latency benchmarks. The full analysis, along with the block size sensitivity table (Table~\ref{tab:block_sensitivity}) and SVD latency scaling table (Table~\ref{tab:latency_comparison}), is detailed in Appendix~\ref{sec:appendix_block_latency}. We establish that block-diagonal SVD with $b=128$ represents the optimal sweet spot for both representation quality and execution speed.
